\section{Models \& Results}
For generel theory on HMM's, see \cite{zucchini2016hidden} and \cite{moler2024statistical} chapter 11. 
All the models implemented are 4 state Gaussian HMMs (emission density is a Gaussian distribution) 
and we have kept the the mean value of state 1 fixed at 400 as this is the baseline $\text{CO}_2$ level in the room 
when no one is present and all windows are closed. This is also shown in the figure \ref{fig:app:co2_data}. 



\subsection{Ordinary HMM model}
\subsubsection{Theory}
The ordinary HMM model is the simplest model implemented in this project and the one most commonly used in the literature.
The transition probabilities are constant and do not depend on time, the observations or any covariates, and so does the emission density.
\newline
For this model, we want to estimate the parameters $\theta = \{\Gamma, \mu, \sigma^2\}$ 
where $\Gamma \in \mathbb{R}^{n \times n}$ is the transition matrix, $\mu \in \mathbb{R}^n$ 
is the mean vector of the Gaussian emission density and $\sigma^2 \in \mathbb{R}^n$ 
is the variance vector of the Gaussian emission density.
Because the markov chain is finite, discrete and homogeneous, 
a unique strictly positive stationary distribution $\delta$ exists 
(chappter 1.3.2 \cite{zucchini2016hidden}) and is given by
$$\delta=\mathbf{1}^T(I-\Gamma+E)^{-1}$$
where $E$ is a matrix of ones. \newline 

To handle the natural bounds on the standard deviation parameters and transition proberbilities, the parameters
were reparameterized as follows:
$$\sigma^2_i = \exp(\eta_i)$$
For the transition matrix, only the $n(n-1)$ off-diagonal entries are
parametrized; the diagonal is used as the reference category with
log-odds fixed to zero. Concretely, the unconstrained parameters are
stored as a matrix $\xi \in \mathbb{R}^{n \times (n-1)}$ that holds the
off-diagonal logits. The transition matrix is then obtained by
exponentiating these logits in the off-diagonal positions of an
identity matrix (so the diagonal is $\exp(0) = 1$) and row-normalizing:
$$
\Gamma_{ij} =
\begin{cases}
\dfrac{\exp(\xi_{ij})}{1 + \sum_{k \neq i} \exp(\xi_{ik})}, & j \neq i, \\[6pt]
\dfrac{1}{1 + \sum_{k \neq i} \exp(\xi_{ik})}, & j = i.
\end{cases}
$$
This is equivalent to a row-wise softmax with the diagonal entry pinned
as the reference, and guarantees $\Gamma_{ij} \geq 0$ and
$\sum_j \Gamma_{ij} = 1$ for any $\xi \in \mathbb{R}^{n \times (n-1)}$.
The stationary distribution $\delta$ is not parametrized directly; it
is recomputed from $\Gamma$ at every optimization step via
$\delta = \mathbf{1}^T(I-\Gamma+E)^{-1}$, so reparametrizing $\Gamma$
through $\xi$ implicitly reparametrizes $\delta$.

The unconstrained parameters $\eta \in \mathbb{R}^n$ and
$\xi \in \mathbb{R}^{n \times (n-1)}$ are the quantities that are
optimized during the fitting process. \newline

The $\mu$ parameters are not bounded, but to avoid label switching for 
the gaussitan distrbitions, the mean parameters are ordered as $\mu_1 < \mu_2 < \mu_3 < \mu_4$ by the reparameterization
$$ \mu_1 = \zeta_1 $$
$$\mu_i = \zeta_1 + \sum_{k=2}^{i} \exp(\zeta_k)$$
where $\zeta \in \mathbb{R}^n$ is the unconstrained parameter vector that is optimized during the fitting process. \newline


\subsubsection{Results}